\chaptercap{CONSIDERAÇÕES FINAIS}

Apresentadas as tecnologias, a modelagem e a implementação do sistema nos capítulos anteriores, este capítulo apresenta as conclusões consolidadas ao longo do desenvolvimento, retomando os objetivos propostos na Seção~\ref{sec:objetivos} e avaliando em que medida foram alcançados. Além disso, destacam-se as principais contribuições técnicas, as restrições identificadas no escopo atual e as perspectivas para evolução contínua da arquitetura.

O objetivo geral deste trabalho consistiu em projetar e implementar um sistema web automatizado para emissão, gestão e validação pública de certificados. O produto de software foi desenvolvido com êxito, resultando em um Produto Mínimo Viável (MVP) funcional, testado e alinhado aos requisitos de negócio. Através do acesso via navegador, a plataforma fornece funcionalidades robustas de emissão, gerenciamento de \textit{templates} visuais, \textit{dashboard} gerencial e um portal público de validação. O disparo automático de e-mails não foi implementado no escopo do MVP; a distribuição dos PDFs ocorre via canais institucionais convencionais (como e-mail manual ou compartilhamento de link), enquanto a geração dos certificados e o portal de validação já atendem ao fluxo principal do sistema. Sob a ótica da engenharia de software, conclui-se que o objetivo central foi plenamente atingido.

Em relação aos objetivos específicos propostos:

\begin{itemize}
    \item O módulo de \textbf{\textit{templates} configuráveis} foi concebido com um editor visual \textit{inline} e renderização em tempo real, permitindo personalização dinâmica de atributos visuais (\ref{rf:upload-templates}, \ref{rf:personalizacao-templates});
    \item O \textbf{mecanismo de autenticação pública}, ancorado em códigos únicos e \textit{QR Codes}, foi integrado ao portal de validação, mitigando riscos de fraudes e falsificações (\ref{rf:portal-validacao}, \ref{rf:validar-certificado});
    \item O \textbf{\textit{dashboard} com indicadores gerenciais} foi implantado, fornecendo telemetria sobre estatísticas de emissão, certificados ativos e uso de \textit{templates} (\ref{rf:estatisticas-emissao});
    \item A \textbf{qualidade do software} foi assegurada através de uma suíte abrangente de \textbf{testes automatizados}, englobando testes unitários (atingindo cobertura de 82,6\% no \textit{frontend} e 96,53\% no \textit{backend}) e \textit{end-to-end}, devidamente integrados ao \textit{pipeline} de integração contínua (Seção~4.3);
    \item O \textbf{ambiente de implantação} foi orquestrado em infraestrutura em nuvem (Render), adotando \textit{pipelines} de Integração e Entrega Contínuas (CI/CD) nativos a partir do GitHub, garantindo \textit{deploys} contínuos e resilientes (Seção~4.1).
\end{itemize}

Devido à adoção da abordagem iterativa e às restrições inerentes ao escopo de um MVP, alguns objetivos específicos foram estrategicamente postergados para iterações futuras. O suporte a importação de dados em lote via arquivos (CSV e Excel) e o serviço de mensageria para envio automático de e-mails não foram incorporados ao \textit{backend} nesta versão. Adicionalmente, os mecanismos rígidos de conformidade com a LGPD e a implementação integral da arquitetura \textit{multi-tenant} requerem validações jurídicas e de infraestrutura, sendo mapeados como evoluções necessárias para o \textit{roadmap} futuro.

\subcap{PRINCIPAIS CONTRIBUICOES}

O sistema desenvolvido oferece uma solução tecnológica moderna e eficiente para o problema identificado, cujas principais contribuições são:

\begin{itemize}
    \item \textbf{Digitalização do fluxo de emissão:} a transição de ferramentas de escritório genéricas para um fluxo automatizado e dedicado reduz a fricção operacional, mitigando falhas humanas e retrabalho;
    \item \textbf{Verificabilidade e Transparência:} a introdução de um portal público, acessível via \textit{QR Code}, descentraliza o processo de auditoria de certificados, provendo verificação em tempo real sem ônus operacional para a instituição emissora;
    \item \textbf{Governança de Dados:} a persistência de registros de emissão em banco de dados relacional com carimbo de tempo e autoria garante rastreabilidade total do ciclo de vida dos certificados;
    \item \textbf{Arquitetura Desacoplada e Escalável:} a segregação entre \textit{frontend} (Vite/React) e \textit{backend} (Next.js), orientada a princípios de arquitetura limpa, promove alta coesão, testabilidade e facilita a evolução independente dos microsserviços;
    \item \textbf{Esteira de CI/CD Automatizada:} a implementação de fluxos de CI/CD garante que todo incremento de código seja automaticamente testado (lint, tipo, unitário, integração e E2E) antes de atingir o ambiente de produção.
\end{itemize}

\subcap{LIMITACOES}

Do ponto de vista arquitetural e funcional, o produto atual apresenta as seguintes limitações, oriundas de \textit{trade-offs} assumidos no escopo do MVP:

\begin{itemize}
    \item O fluxo de \textbf{emissão em lote} é atualmente sequencial e manual na interface; a ausência de um \textit{message broker} (processamento em fila) e de suporte à importação de arquivos (CSV/Excel) limita o rendimento para volumes massivos de dados;
    \item A \textbf{notificação ativa (E-mail)} não foi implementada no lado do servidor (sem integração com provedores SMTP). O artefato atual no \textit{frontend} (\code{envios.ts}) apenas estrutura o agendamento em armazenamento local (\code{localStorage}) de forma simulada;
    \item A \textbf{arquitetura \textit{multi-tenant}} não foi finalizada — o banco de dados opera sob o paradigma de \textit{tenant} único, sem o particionamento lógico de dados essencial para um modelo SaaS em larga escala;
    \item A \textbf{adequação à LGPD} carece de refinamento — mecanismos como ofuscação/anonimização de dados pessoais e gestão granular de consentimento ainda não estão materializados no código;
    \item A \textbf{reemissão com manutenção de histórico} (\ref{rf:reemissao-historico}) não foi implementada — a edição de certificados sobrescreve os dados in-place, sem gerar novo código de verificação ou preservar a versão anterior;
    \item A \textbf{autenticação via JWT com Supabase Auth} não foi integralmente implementada — o middleware de autenticação existe no backend e foi validado em testes unitários, mas o fluxo completo de login, gerenciamento de sessão e renovação de tokens está previsto para versões futuras;
    \item A suíte de testes carece de \textbf{testes de carga e estresse}, os quais são fundamentais para atestar o comportamento e as garantias do serviço de geração de PDFs sob alta concorrência.
\end{itemize}

\subcap{TRABALHOS FUTUROS}

Vislumbrando a evolução contínua da plataforma, recomenda-se a exploração dos seguintes aprimoramentos técnicos e funcionais:

\begin{itemize}
    \item Materialização das entidades \textbf{PARTICIPANTES} e \textbf{CURSOS} no banco de dados, substituindo a abordagem denormalizada atual (campos \code{recipientName}, \code{recipientCPF}, \code{courseName} e \code{courseHours} diretamente na tabela \code{CERTIFICADOS}) por um modelo relacional normalizado com chaves estrangeiras (\code{participantId} e \code{courseId}), habilitando histórico de certificados por participante, listagem por curso e relatórios gerenciais por unidade de treinamento;
    \item Adoção de um \textbf{serviço de mensageria / fila} (como Redis/BullMQ ou RabbitMQ) para orquestrar o processamento assíncrono de emissões em lote e processamento de arquivos CSV/Excel;
    \item Integração com provedores de nuvem para \textbf{envio transacional de e-mails}, implementando retentativas exponenciais (\textit{exponential backoff}) e \textit{templates} ricos;
    \item Refatoração do modelo de dados para suportar plenamente o \textbf{isolamento \textit{multi-tenant}}, com políticas de segurança de linha (\textit{Row Level Security} - RLS) via Supabase;
    \item Implementação nativa de um módulo de \textbf{auditoria e conformidade com a LGPD}, incorporando \textit{logs} imutáveis e fluxos automatizados para a exclusão de dados pessoais;
    \item Execução de campanhas de \textbf{testes não-funcionais} (carga/desempenho) e \textbf{análise automatizada de vulnerabilidades} na esteira de CI/CD, com ênfase em patamares de concorrência superiores (50, 100 e 500 requisições simultâneas) para o serviço de geração de PDF, validando o comportamento do sistema sob carga real de produção;
    \item Investigação sobre a viabilidade de ancoragem criptográfica dos certificados através de \textbf{Assinatura Digital (ICP-Brasil)} ou redes \textbf{Blockchain}, visando segurança e imutabilidade incontestáveis;
    \item Implementação de \textbf{versionamento da API} (ex.: \texttt{/api/v1/}, \texttt{/api/v2/}), permitindo evolução sem quebrar contratos com clientes existentes e habilitando migração gradual entre versões.
\end{itemize}
